\section{Objetivos}

\subsection{Objetivo Geral}

Consolidar os fundamentos teóricos e computacionais da computação
quântica aplicada à simulação molecular, por meio da compreensão,
implementação e análise do algoritmo \textit{Variational Quantum
Eigensolver} (VQE), estabelecendo uma prova de conceito para
simulações quânticas híbridas aplicadas a sistemas moleculares
modelo.

\subsection{Objetivos Específicos}

\begin{itemize}
    \item Consolidar a fundamentação teórica em computação quântica aplicada à simulação de sistemas físicos, incluindo estados quânticos, circuitos e princípio variacional;
    
    \item Formular o problema de estrutura eletrônica molecular como um problema de autovalor do Hamiltoniano eletrônico;
    
    \item Estudar e analisar o algoritmo variacional VQE no contexto NISQ, compreendendo seu funcionamento e limitações;
    
    \item Compreender e estruturar o fluxo híbrido quântico-clássico do VQE: construção do Hamiltoniano, escolha do ansatz, medição de observáveis e otimização clássica;
    
    \item Implementar uma prova de conceito do VQE no Qiskit para moléculas modelo (H$_2$ e LiH);
    
    \item Desenvolver domínio prático do ecossistema IBM Quantum para execução local e em hardware real;
    
    \item Estabelecer bases metodológicas para extensão futura a sistemas moleculares mais complexos e integração com técnicas de aprendizado de máquina;
    
    \item Garantir reprodutibilidade científica por meio de versionamento, documentação e disponibilização pública dos códigos desenvolvidos.
\end{itemize}
